\section{Introducción}

\subsection{Propósito}
El propósito del presente documento es definir de manera formal y detallada la \textbf{Especificación de Requisitos de Software (SRS)} para el sistema \textbf{``Gemelo Digital EAM''}, como parte de la fase de \textbf{Análisis de Modelo Conceptual}. Este documento integra el refinamiento de 30 Historias de Usuario y la derivación de 360 requisitos técnicos alineados con estándares internacionales de gestión de activos y seguridad industrial.

Este SRS constituye la línea base técnica para la transición hacia la fase de implementación, asegurando que la solución trascienda la digitalización de procesos y establezca un marco de gobernanza basado en datos para el mantenimiento de activos críticos en el sector industrial.

\subsection{Ámbito del Sistema (Alcance)}
El sistema \textbf{``Gemelo Digital EAM''} es una plataforma híbrida diseñada para optimizar la gestión operativa y estratégica mediante la convergencia de datos de sensores (IoT), modelos 3D y motores de inferencia IA.

Esta solución busca transformar los procesos manuales o desconectados que actualmente impactan la eficiencia de la planta, específicamente:

\begin{itemize}
    \item \textbf{Captura de Datos en Campo:} Reemplazo de checklists en papel por formularios digitales que alimentan indicadores en tiempo real.
    \item \textbf{Ingesta de Informes Externos:} Automatización de la extracción de datos de reportes de contratistas mediante inteligencia artificial.
    \item \textbf{Acceso a Información Técnica:} Disponibilidad inmediata de planos y manuales en el punto de trabajo vía dispositivos móviles.
    \item \textbf{Verificación de Stock:} Validación automática de repuestos antes de la programación de intervenciones.
\end{itemize}
\textbf{El sistema SÍ incluirá:}

\begin{enumerate}
    \item \textbf{Gobernanza de Mantenimiento:} Ciclo de vida completo de órdenes de trabajo con priorización objetiva mediante matriz RIME y cierre digital bajo taxonomía ISO 14224.
    \item \textbf{Control de Recursos:} Gestión de activos serializados (Equipment Units) y control de stock crítico basado en clasificación ABC y prevención de duplicados vía Commodity Codes.
    \item \textbf{Visualización y Navegación:} Gemelo digital interactivo con navegación espacial por zonas funcionales, incluyendo capas dinámicas de seguridad para permisos de trabajo (PTW) y rutas de aislamiento (LOTO) con lógica de bloqueo activo.
    \item \textbf{Monitoreo de Condición:} Lectura y visualización de variables de telemetría provenientes de sensores IoT para la detección de anomalías y alimentación del mantenimiento predictivo.
    \item \textbf{Inteligencia y Auditoría:} Consola de administración con control de acceso RBAC granular, auditoría inmutable con cadena de integridad criptográfica (Hash-Chaining) y dashboards de KPIs (MTBF, MTTR).
    \item \textbf{Framework de Asistencia IA:} Motor de asistencia técnica basado en RAG (Retrieval-Augmented Generation) para consulta de manuales en lenguaje natural.
\end{enumerate}
\textbf{El sistema NO incluirá (Exclusiones y Consideraciones Futuras):}

\begin{enumerate}
    \item \textbf{Edición de Malla 3D:} El sistema consume modelos existentes pero no permite el modelado geométrico interno.
    \item \textbf{Integración ERP Financiera:} La sincronización automática con sistemas contables externos se difiere a fases posteriores debido a dependencias de APIs de terceros.
    \item \textbf{Control Industrial Activo (SCADA/HMI):} Capacidad para enviar comandos de parada o arranque directo a la maquinaria. Se excluye explícitamente en esta fase para garantizar la seguridad funcional de la planta, operando en modo de ``monitoreo''.
    \item \textbf{Integración con Suites de Productividad:} Sincronización nativa con plataformas como Google Workspace u Office 365 para la gestión de identidad y notificaciones avanzadas (Diferido a versiones posteriores).
\end{enumerate}
\subsection{Definiciones, Acrónimos y Abreviaturas}

\textbf{Nomenclatura de Módulos}

\begin{itemize}
    \item \textbf{ADM:} Administración e Inteligencia de Negocio.
    \item \textbf{INV:} Gestión de Recursos e Inventario.
    \item \textbf{MTTO:} Gestión de Operaciones de Mantenimiento.
    \item \textbf{VIS:} Visualización y Gemelo Digital.
    \item \textbf{TR:} Requisito Transversal (Arquitectura Común).
\end{itemize}
\textbf{Términos Técnicos y del Dominio}

\begin{itemize}
    \item \textbf{API (Application Programming Interface):} Interfaz que permite la comunicación e intercambio de datos entre dos sistemas de software independientes.
    \item \textbf{CMMS (Computerized Maintenance Management System):} Software diseñado para simplificar la gestión del mantenimiento de activos.
    \item \textbf{DT (Digital Twin / Gemelo Digital):} Representación virtual dinámica de un activo físico o sistema que se actualiza con datos en tiempo real.
    \item \textbf{EAM (Enterprise Asset Management):} Gestión de activos empresariales; software que gestiona el ciclo de vida completo de los activos físicos.
    \item \textbf{ERP (Enterprise Resource Planning):} Sistema de planificación de recursos empresariales utilizado para la gestión de procesos de negocio (finanzas, compras).
    \item \textbf{HSEQ:} Acrónimo para Salud, Seguridad, Medio Ambiente y Calidad (Health, Safety, Environment, Quality).
    \item \textbf{KPI (Key Performance Indicator):} Indicador clave de rendimiento utilizado para medir el éxito de una actividad.
    \item \textbf{LOTO (Lockout/Tagout):} Procedimiento de seguridad para asegurar que las máquinas estén debidamente apagadas y no se puedan encender durante el mantenimiento.
    \item \textbf{MRO:} Mantenimiento, Reparación y Operaciones; se refiere a los suministros y repuestos necesarios para mantener la planta operativa.
    \item \textbf{MTBF (Mean Time Between Failures):} Tiempo medio entre fallas; métrica de confiabilidad que mide el tiempo promedio que un equipo funciona sin fallar.
    \item \textbf{MTTR (Mean Time To Repair):} Tiempo medio de reparación; métrica que mide el tiempo promedio necesario para reparar un activo.
    \item \textbf{OEE (Overall Equipment Effectiveness):} Efectividad Global del Equipo; estándar para medir la productividad de la manufactura.
    \item \textbf{OT / WO (Orden de Trabajo / Work Order):} Documento que autoriza y detalla una tarea de mantenimiento específica.
    \item \textbf{PTW (Permit to Work):} Permiso de Trabajo; autorización formal para realizar actividades de alto riesgo.
    \item \textbf{RAG (Retrieval-Augmented Generation):} Técnica de IA que optimiza las respuestas de un LLM consultando una base de conocimientos específica.
    \item \textbf{SIL-2 (Safety Integrity Level 2):} Nivel de integridad de seguridad aplicado a la lógica de software crítica.
    \item \textbf{Sidecar Pattern:} Patrón de diseño que desacopla la visualización de la inteligencia de datos.
\end{itemize}
\subsection{Referencias}
\begin{sloppypar}

\textbf{Normas y Estándares Técnicos}

\begin{itemize}
    \item International Organization for Standardization. (2023). \textit{Systems and software engineering — Systems and software Quality Requirements and Evaluation (SQuaRE) — Product quality model} (ISO/IEC 25010:2023).
    \item International Organization for Standardization. (2016). \textit{Petroleum, petrochemical and natural gas industries — Collection and exchange of reliability and maintenance data for equipment} (ISO 14224:2016).
    \item International Organization for Standardization. (2014). \textit{Asset management — Management systems — Requirements} (ISO 55001:2014).
    \item International Electrotechnical Commission. (2010). \textit{Functional safety of electrical\slash electronic\slash programmable electronic safety-related systems} (IEC 61508).
\end{itemize}
\textbf{Bibliografía Académica}

\begin{itemize}
    \item Chen, S., Turanoglu Bekar, E., Bokrantz, J., \& Skoogh, A. (2025). AI-enhanced digital twins in maintenance: Systematic review, industrial challenges, and bridging research–practice gaps. \textit{Journal of Manufacturing Systems}, 82, 678-699. \url{https://doi.org/10.1016/j.jmsy.2025.07.006}
    \item Cohn, M. (2004). \textit{User Stories Applied: For Agile Software Development}. Addison-Wesley Professional.
    \item Hollifield, B., Oliver, D., Nimmo, I., \& Habibi, E. (2008). \textit{The high performance HMI handbook: A comprehensive guide to designing, implementing and maintaining effective HMIs for industrial plant operations}. Plant Automation Services, Inc.
    \item Wireman, T. (2004). \textit{Benchmarking best practices in maintenance management}. Industrial Press Inc.
    \item Wireman, T. (2008). \textit{MRO inventory and purchasing}. Industrial Press Inc.
\end{itemize}
\end{sloppypar}
